% ======================================================================
% 第 2 章：标准模型中的粒子相互作用
% 原著：R. Gupta, "Introduction to Lattice QCD" (arXiv:hep-lat/9807028)，
% 第 2 节，第 7–9 页。
% ======================================================================

\chapter{标准模型中的粒子相互作用}

粒子相互作用的标准模型（SM）是自然界四种力中三种的综合。这些力由规范理论描述，
每一种都由一个耦合常数来刻画，如下所列。

\begin{center}
\begin{tabular}{ll}
强相互作用 & $\alpha_s \sim 1$ \\
电磁相互作用 & $\alpha_{em} \approx 1/137$ \\
弱相互作用 & $G_F \approx 10^{-5}\,\mathrm{GeV}^{-2}$.
\end{tabular}
\end{center}

物质的基本组分是六种夸克 $u$、$d$、$s$、$c$、$b$、$t$，每种夸克都有 3 种颜色，
以及六种轻子 $e$、$\nu_e$、$\mu$、$\nu_\mu$、$\tau$、$\nu_\tau$。夸克和轻子被分为
3 代族。这些粒子之间的相互作用由矢量玻色子来传递：8 种胶子传递强相互作用，
$W^\pm$ 和 $Z$ 传递弱相互作用，而电磁相互作用由光子 $\gamma$ 承载。弱玻色子通过
希格斯机制获得质量，并且在最小 SM 表述中应当存在一个中性希格斯粒子。这一点迄今
尚未在实验上观察到。描述这些组分之间相互作用的数学结构是一种定域规范场论，其规范群
为 $SU(3) \times SU(2) \times U(1)$。刻画 SM 的参数如下所列；它们的起源在于一个
更为基本但尚未被发现的理论。

\begin{table}[htbp]
\centering
\caption*{标准模型中的参数}
\label{tab:sm-parameters}
\begin{tabular}{l c p{7cm}}
\toprule
参数 & 数目 & 说明 \\
\midrule
夸克质量 & 6 & \begin{tabular}[t]{@{}l@{}}
$u$, $d$, $s$ 轻 \\ $c$, $b$ 重 \\ $t = 175 \pm 6$~GeV \end{tabular} \\
轻子质量 & 6 & \begin{tabular}[t]{@{}l@{}}
$e$, $\mu$, $\tau$ \\ $M_{\nu_e}$, $\nu_\mu$, $\nu_\tau = ?$ \end{tabular} \\
$W^\pm$ 的质量 & 1 & 80.3~GeV \\
$Z$ 的质量 & 1 & 91.2~GeV \\
胶子、$\gamma$ 的质量 & & 0（规范对称性） \\
希格斯质量 & 1 & 尚未观测到 \\
耦合常数 $\alpha_s$ & 1 & 能量 $\lesssim 1$~GeV 时 $\approx 1$ \\
耦合常数 $\alpha_{em}$ & 1 & 1/137（在 $M_Z$ 处 = 1/128.9） \\
耦合常数 $G_F$ & 1 & $10^{-5}\,\mathrm{GeV}^{-2}$ \\
弱混合角 & 3 & $\theta_1$, $\theta_2$, $\theta_3$ \\
CP 破坏相位 & 1 & $\delta$ \\
强 CP 参数 & 1 & $\Theta = ?$ \\
\bottomrule
\end{tabular}
\end{table}

问号表示这些值极其微小且尚未被测量。轻子扇区没有列出混合角，因为我假定中微子质量
均为零。关于中微子质量及其混合的实验搜索现状，本学校的 Prof.~Sciulli 已作了报告。
弱相互作用的结构以及弱混合角的起源，由 Daniel Treille 作了综述，而 David Kosower
则讲解了微扰 QCD。

在 SM 中，$W$ 玻色子与夸克的带电流相互作用由下式给出
\begin{equation}
H_W = \frac{1}{2} \left( J_\mu W^\mu + \text{h.c.} \right)
\label{eq:hw}
\end{equation}
其中流为
\begin{equation}
J_\mu = \frac{g_w}{2\sqrt{2}} \left( \bar{u},\ \bar{c},\ \bar{t} \right)
\gamma_\mu (1 - \gamma_5)\, V
\begin{pmatrix} d \\ s \\ b \end{pmatrix}
\label{eq:jmu}
\end{equation}
而费米耦合常数与 $SU(2)$ 耦合的关系为 $G_F / \sqrt{2} = g_w^2 / 8 M_W^2$。
$V$ 是 $3 \times 3$ 的 Cabibbo-Kobayashi-Maskawa（CKM）矩阵，它在味变换矩阵中
有一个简单的表示
\begin{equation}
V_{CKM} =
\begin{pmatrix}
V_{ud} & V_{us} & V_{ub} \\
V_{cd} & V_{cs} & V_{cb} \\
V_{td} & V_{ts} & V_{tb}
\end{pmatrix}
\label{eq:ckm}
\end{equation}
这个矩阵具有非对角元，因为弱相互作用的本征态与强相互作用的本征态并不相同，即
$d$、$s$、$b$ 夸克（或等价地 $u$、$c$、$t$ 夸克）在弱相互作用下发生混合。

对于 3 代夸克，幺正条件 $V^{-1} = V^\dagger$ 施加了约束。因此只有四个独立参数，
可以用 4 个角 $\theta_1$、$\theta_2$、$\theta_3$ 以及 CP 破坏相位 $\delta$ 来表示。
最初由 Kobayashi 和 Masakawa 提出的表示是
\begin{equation}
V_{CKM} =
\begin{pmatrix}
c_1 & -s_1 c_3 & -s_1 s_3 \\
s_1 c_2 & c_1 c_2 c_3 - s_2 s_3 e^{i\delta} & c_1 c_2 s_3 + s_2 c_3 e^{i\delta} \\
s_1 s_2 & c_1 s_2 c_3 + c_2 s_3 e^{i\delta} & c_1 s_2 s_3 - c_2 c_3 e^{i\delta}
\end{pmatrix}
\label{eq:km}
\end{equation}
其中对 $i = 1, 2, 3$ 有 $c_i = \cos\theta_i$、$s_i = \sin\theta_i$。这个矩阵在
现象学上更有用的表示是 Wolfenstein 参数化（精确到 $O(\lambda^4)$）
\begin{equation}
V_{CKM} =
\begin{pmatrix}
1 - \lambda^2/2 & \lambda & \lambda^3 A (\rho - i\eta) \\
-\lambda & 1 - \lambda^2/2 & \lambda^2 A \\
\lambda^3 A (1 - \rho - i\eta) & -\lambda^2 A & 1
\end{pmatrix}
+ O(\lambda^4)
\label{eq:wolfenstein}
\end{equation}
其中 $\lambda = \sin\theta_c = 0.22$ 已知到 1\% 的精度，$A \approx 0.83$ 仅已知到
10\% 的精度，而 $\rho$ 和 $\eta$ 知之甚少。式~\eqref{eq:wolfenstein} 中给出的
CKM 矩阵元 $V_{td}$ 和 $V_{ub}$ 在 $\eta \neq 0$ 时是复数，因此 SM 中存在着一种
自然的 CP 破坏机制。CKM 矩阵的现象学以及这四个参数的测定，本学校的
Profs.~Buras 和 Richman 已作了综述 [10,11]。

图~\ref{fig:1} 显示了那些对 CKM 参数最为敏感的探测过程的一览。格点QCD 计算正开始
为与这些过程相关的强子矩阵元提供最可靠的估计之一。这些计算的细节及其对确定 $\rho$
和 $\eta$ 参数的影响，将由本学校的 Prof.~Martinelli 讲解。我的目标是为你提供
必要的背景知识。
